
\documentclass[10pt,
% a4paper,
%twocolumn,
fleqn,
%landscape, 
% papersize,
dvipdfmx,
uplatex
]{jsarticle}



\def\maru#1{\textcircled{\scriptsize#1}}%丸囲み番号

% \RequirePackage[2020/09/30]{platexrelease}

%太字設定
\usepackage[deluxe]{otf}

\usepackage{emathEy}

\usepackage[g]{esvect}

%大きな文字
\usepackage{fix-cm}

%定理環境
\usepackage{emathThm}
%\theoremstyle{boxed}
\theorembodyfont{\normalfont}
\newtheorem{Question}{問題}[subsection]
\newtheorem{Q}{}[subsection]
\newtheorem{question}[Question]{}
\newtheorem{quuestion}{}[subsection]

%セクション，大問番号のデザイン
\renewcommand{\labelenumi}{(\arabic{enumi})}
\renewcommand{\theenumii}{\alph{enumii})}
\renewcommand{\thesection}{第\arabic{section}章}

%用紙サイズの詳細設定
\usepackage{bxpapersize}
\papersizesetup{size={80mm,45mm}}
\usepackage[top=0.7zw,bottom=0truemm,left=3truemm,right=133truemm]{geometry}
\usepackage[dvipdfmx]{graphicx}

%余白など
\usepackage{setspace} % 行間
\setlength{\mathindent}{1zw}
\setlength\parindent{0pt}


%色カラーに関する設定
\usepackage{color}
\definecolor{shiro}{rgb}{0.95703125,0.87109375,0.7421875}
\definecolor{kin}{rgb}{0.95703125,0.87109375,0.7421875}
\definecolor{orange}{rgb}{1,0.7,0.2}
\definecolor{bradorange}{rgb}{1,0.5,0}
\definecolor{pink}{rgb}{0.9176,0.5686,0.5960}
\definecolor{mizu}{rgb}{0.6156,0.8,0.9955}
\color{kin}
% \pagecolor{hukamido}

\usepackage{at}%図の配置
% \usepackage{wallpaper}

\begin{document}

\bf\boldmath



\at(0cm,0cm){\includegraphics[width=8cm,bb=0 0 1920 1080]{./youtube/thumbnails/templates/smart_background/複素数と方程式.jpeg}}
\at(6.4cm,0.2cm){\small\color{bradorange}$\overset{\text{複素数と方程式}}{\text{応用}}$}
{\color{orange}\LARGE\underline{$3$元対称式の連立方程式}}\vspace{0.3zw}

\large
問．

\Large
\vspace{-1.1zw}
\hspace{1zw}$\left\{\begin{array}{l}x+y+z=-a-2,\;\vspace{-0.2zw}\\ xy+yz+zx=2a+1,\;\vspace{-0.2zw}\\ xyz=-2\end{array}\right.\vspace{0.0zw}$

\large
が，$x,\;y,\;z$が実数である解をもつ\\
\hfill ような実数$a$の範囲を求めよ．


\newpage



\at(0cm,0cm){\includegraphics[width=8cm,bb=0 0 1920 1080]{./youtube/thumbnails/templates/smart_background/複素数と方程式.jpeg}}
\at(6.4cm,0.2cm){\small\color{bradorange}$\overset{\text{複素数と方程式}}{\text{応用}}$}
{\color{orange}\LARGE\underline{$3$解から$3$次方程式を作る}}\vspace{0.3zw}

\large
問．方程式$x^3-x+1=0$の\\
\hfill $3$つの解を$\alpha ,\;\beta ,\;\gamma$とするとき，

\LARGE 
\vspace{0.2zw}
\hspace{0.5zw}$\bunsuu{1}{\alpha },\;\bunsuu{1}{\beta },\;\bunsuu{1}{\gamma}$を$3$解にもつ\vspace{0.2zw}

\large
\hfill ような$3$次方程式を$1$つ作れ．


\newpage



\at(0cm,0cm){\includegraphics[width=8cm,bb=0 0 1920 1080]{./youtube/thumbnails/templates/smart_background/複素数と方程式.jpeg}}
\at(6.4cm,0.2cm){\small\color{bradorange}$\overset{\text{複素数と方程式}}{\text{応用}}$}
{\color{orange}\large\underline{$3$次方程式の$3$解$\alpha ,\;\beta ,\;\gamma$の対称式}}\vspace{0.3zw}

\large
問．方程式$x^3-3x+1=0$の\\
\hfill $3$つの解を$\alpha ,\;\beta ,\;\gamma$とするとき，

\huge
\vspace{0.3zw}
\hspace{0.1zw}
$\begin{array}{l}\alpha ^3+\beta ^3+\gamma^3\;\\
\alpha ^4+\beta ^4+\gamma^4
\end{array}$

\large
\vspace{-1.5zw}
\hfill を求めよ．


\newpage



\at(0cm,0cm){\includegraphics[width=8cm,bb=0 0 1920 1080]{./youtube/thumbnails/templates/smart_background/複素数と方程式.jpeg}}
\at(6.4cm,0.2cm){\small\color{bradorange}$\overset{\text{複素数と方程式}}{\text{応用}}$}
\at(2cm,0.2cm){\color{orange}\LARGE\underline{}}\vspace{0.3zw}

\scriptsize
\hspace{0zw}$x^2 = 1$\vspace{-0.5zw}

\large 
\hspace{0.5zw}$x^3 = 1$\vspace{-0.5zw}

\LARGE 
\hspace{1.2zw}$x^4 = 1$\vspace{-0.5zw}

\HUGE
\hspace{1.8zw}{\color{orange} $x^5 = 1$}\vspace{0.1zw}

\fontsize{45}{10} \selectfont
\hspace{2.6zw}$x^6 =$



\newpage



\at(0cm,0cm){\includegraphics[width=8cm,bb=0 0 1920 1080]{./youtube/thumbnails/templates/smart_background/複素数と方程式.jpeg}}
\at(6.4cm,0.2cm){\small\color{bradorange}$\overset{\text{複素数と方程式}}{\text{応用}}$}
{\color{orange}\huge\underline{相反方程式}}\vspace{0.1zw}

\small 
問．次の$4$次方程式について

\Large
\vspace{0.1zw}
\hspace{0.5zw}$x^4-4x^3+7x^2-8x+4=0\vspace{0.5zw}$

\small 
(1)  両辺を$x^2$で割って，$t=x+\bunsuu{2}{x}$とおくことで\\
\hfill 得られる$t$に関する$2$次方程式を求めよ．\\
(2)  解をすべて求めよ．\\



\newpage



\at(0cm,0cm){\includegraphics[width=8cm,bb=0 0 1920 1080]{./youtube/thumbnails/templates/smart_background/複素数と方程式.jpeg}}
\at(6.4cm,0.2cm){\small\color{bradorange}$\overset{\text{複素数と方程式}}{\text{応用}}$}
{\color{orange}\huge\underline{複$2$次方程式}}\vspace{0.3zw}

\Large
問．方程式

\Huge
\vspace{0.3zw}
\hspace{0.2zw}$z^4-6z^2+{25}=0\vspace{0.3zw}$

\Large
\hfill を解け．


\newpage



\at(0cm,0cm){\includegraphics[width=8cm,bb=0 0 1920 1080]{./youtube/thumbnails/templates/smart_background/複素数と方程式.jpeg}}
\at(6.4cm,0.2cm){\small\color{bradorange}$\overset{\text{複素数と方程式}}{\text{応用}}$}
{\color{orange}\LARGE\underline{複素数係数の$2$次方程式}}\vspace{0.3zw}

\normalsize 
問．(1)  $z=x+yi(x,\;y$は実数$)$が，\\
\hfill $z^2=i$を満たすように，$x,\;y$の値を定めよ．\\
(2)  $2$次方程式

\LARGE
\vspace{0.0zw}
\hspace{0.5zw}$w^2+2\left(1+i\right)w+i=0\vspace{0.1zw}$

\normalsize 
\hfill を解け．\\



\newpage



\at(0cm,0cm){\includegraphics[width=8cm,bb=0 0 1920 1080]{./youtube/thumbnails/templates/smart_background/複素数と方程式.jpeg}}
\at(6.4cm,0.2cm){\small\color{bradorange}$\overset{\text{複素数と方程式}}{\text{応用}}$}
{\color{orange}\Large\underline{複素数係数の方程式の実数解}}\vspace{0.3zw}

\normalsize
問．$a$を実数の定数とする．$x$の$2$次方程式
 
\vspace{0.1zw}
\hspace{0.5zw}$\left(1+i\right)x^2-\left(a+1+i\right)x+\left(2-ai\right)=0\vspace{0zw}$

\LARGE
が実数解をもつのは\vspace{-0.2zw}\\
\hfill $a$がどんな値のときか．

\normalsize
\hfill ただし，$i$は虚数単位である．


\newpage



\at(0cm,0cm){\includegraphics[width=8cm,bb=0 0 1920 1080]{./youtube/thumbnails/templates/smart_background/複素数と方程式.jpeg}}
\at(6.4cm,0.2cm){\small\color{bradorange}$\overset{\text{複素数と方程式}}{\text{応用}}$}
{\color{orange}\Large\underline{整式の割り算の余り Lv.2 }}\vspace{0.3zw}

\small
問．整式$P\left(x\right)$を$x-1$で割ったときの余りが$5$,\;\\
\hfill $\left(x+1\right)^2$で割ったときの余りが$x-8$であるとき，

\huge 
\vspace{-0.1zw}
$\left(x-1\right)\left(x+1\right)^2$で\vspace{-0.1zw}\\
\hfill 割ったときの余り
\vspace{0.1zw}

\small
\hfill を求めよ．


\newpage



\at(0cm,0cm){\includegraphics[width=8cm,bb=0 0 1920 1080]{./youtube/thumbnails/templates/smart_background/複素数と方程式.jpeg}}
\at(6.4cm,0.2cm){\small\color{bradorange}$\overset{\text{複素数と方程式}}{\text{応用}}$}
{\color{orange}\Large\underline{重解をもつ$3$次方程式 Lv.2 }}\vspace{0.1zw}

\Large
問．$3$次方程式

\Huge
\vspace{0.0zw}
\hspace{.5zw}$x^3-2x+k=0\vspace{0.1zw}$

\Large
が重解をもつのは，\\
\hfill $k$がいかなる値のときか．


\newpage



\at(0cm,0cm){\includegraphics[width=8cm,bb=0 0 1920 1080]{./youtube/thumbnails/templates/smart_background/複素数と方程式.jpeg}}
\at(6.4cm,0.2cm){\small\color{bradorange}$\overset{\text{複素数と方程式}}{\text{応用}}$}
{\color{orange}\huge\underline{$3$次方程式の虚数解}}\vspace{0.3zw}

\normalsize 
問．$a,\;b$を実数とする．方程式

\LARGE
\vspace{0.1zw}
\hspace{0.5zw}$2x^3+ax^2+bx-6=0\vspace{0.1zw}$
\hfill 

\Large
が$x=1+i$を解にもつとき，\\
\hfill 残りの解を求めよ．


\newpage



\at(0cm,0cm){\includegraphics[width=8cm,bb=0 0 1920 1080]{./youtube/thumbnails/templates/smart_background/複素数と方程式.jpeg}}
\at(6.4cm,0.2cm){\small\color{bradorange}$\overset{\text{複素数と方程式}}{\text{応用}}$}
{\color{orange}\Large\underline{$3$次方程式の解と係数の関係}}\vspace{0.3zw}

\small 
問．$3$次方程式$ax^3+bx^2+cx+d=0$の$3$解が$x=\alpha ,\;\beta ,\;\gamma$であること，次の条件は同値である．

\large
\vspace{0.3zw}
\hspace{1.5zw}$\left\{\begin{array}{l} \alpha +\beta +\gamma=-\bunsuu{b}{a},\;\vspace{-0.2zw}\\ \alpha \beta +\beta \gamma+\gamma\alpha =\bunsuu{c}{a},\;\vspace{-0.2zw}\\ \alpha \beta \gamma=-\bunsuu{d}{a}\end{array}\right.\vspace{0.3zw}$




\end{document}

